\section{Figurer og referencer}

For at indsætte en figur benyttes følgende kode
\begin{figure}[H]
    \centering
    \includegraphics[scale=0.3]{Billeder/DTULogo.png}
    \caption{DTU logo}
    \label{fig:DTU_logo}
\end{figure}
Hvis man ønsker at lave en reference til en figur eller graf, kan man bruge dets label således: \autoref{fig:DTU_logo} eller \ref{fig:DTU_logo}.

Man kan også indsætte flere billeder sammen: 
\begin{figure}[H] % når man tilføjer H, betyder det at man ønkser figuren bliver placeret lige her - ellers har billeder en tendens til at flyve rundt i et latex dokument
     \centering
     \begin{subfigure}[b]{0.3\textwidth}
         \centering
         \includegraphics[width=\textwidth]{Billeder/abe.jpeg}
         \caption{Klog abe}
         \label{fig:abe}
     \end{subfigure}
     \hfill
     \begin{subfigure}[b]{0.3\textwidth}
         \centering
         \includegraphics[width=\textwidth]{Billeder/kanin.jpeg}
         \caption{Cute kanin}
         \label{fig:kanin}
     \end{subfigure}
     \hfill
     \begin{subfigure}[b]{0.3\textwidth}
         \centering
         \includegraphics[width=\textwidth]{Billeder/kat.jpeg}
         \caption{En lidt sad kat}
         \label{fig:kat}
     \end{subfigure}
        \caption{Tre meget imponerende dyr}
        \label{fig:three graphs}
\end{figure}
Aben findes her: \ref{fig:abe}